%==================================================================
% BAB 2 TINJAUAN PUSTAKA
% Berisi teori pendukung dan penelitian terdahulu.
% Perhatikan contoh penulisan rumus, tabel, dan sitasi di bab ini.
%==================================================================
\chapter{TINJAUAN PUSTAKA}
\label{chap:tinjauanpustaka}

%------------------------------------------------------------------
\section{Kesegaran Ikan}
%------------------------------------------------------------------

Kesegaran ikan adalah kondisi mutu ikan yang ditentukan oleh perubahan
biokimia dan mikrobiologis setelah ikan ditangkap. Menurut
\citet{estiasih2015}, kemunduran mutu ikan dipengaruhi oleh suhu
penyimpanan, cara penanganan, dan kandungan lemak pada tubuh ikan.
Semakin tinggi suhu penyimpanan, semakin cepat pula proses kemunduran
mutu berlangsung.

Secara organoleptik, kesegaran ikan dapat dinilai melalui penampakan
mata, warna insang, tekstur daging, dan bau. Ikan segar memiliki mata
yang jernih dan cembung, insang berwarna merah cerah, serta daging yang
kenyal. Seiring bertambahnya waktu penyimpanan, mata menjadi keruh dan
cekung, insang berubah pucat hingga kecoklatan, dan daging menjadi
lembek. Perubahan warna tersebut menjadi dasar pemilihan ciri warna
pada penelitian ini.

%------------------------------------------------------------------
\section{Pengolahan Citra Digital}
%------------------------------------------------------------------

Pengolahan citra digital merupakan proses mengolah citra menjadi citra
lain atau menjadi data yang dapat diinterpretasikan oleh komputer
\citep{gonzalez2018}. Citra digital direpresentasikan sebagai fungsi
dua dimensi $f(x,y)$ dengan $x$ dan $y$ menyatakan koordinat spasial
dan nilai fungsi menyatakan intensitas atau tingkat kecerahan piksel
pada koordinat tersebut.

Tahapan pengolahan citra yang digunakan pada penelitian ini meliputi
penyekalaan ukuran citra, perbaikan kontras, pemisahan kanal warna, dan
ekstraksi ciri statistik. Penyekalaan ukuran dilakukan agar seluruh
citra masukan memiliki dimensi yang seragam sebelum diproses lebih
lanjut.

%------------------------------------------------------------------
\section{Ruang Warna HSV}
%------------------------------------------------------------------

Ruang warna HSV memisahkan informasi warna (\textit{hue}) dari
informasi kecerahan (\textit{value}), sehingga lebih tahan terhadap
perubahan intensitas cahaya dibandingkan ruang warna RGB. Konversi
nilai RGB ke HSV dihitung menggunakan \cref{eq:hsv}.

\begin{equation}
\begin{aligned}
V &= \max(R,G,B) \\
S &= \begin{cases}
       \dfrac{V-\min(R,G,B)}{V}, & V \neq 0 \\[6pt]
       0, & V = 0
     \end{cases} \\
H &= \begin{cases}
       60^{\circ}\left(\dfrac{G-B}{V-\min(R,G,B)}\right), & V = R \\[6pt]
       60^{\circ}\left(2+\dfrac{B-R}{V-\min(R,G,B)}\right), & V = G \\[6pt]
       60^{\circ}\left(4+\dfrac{R-G}{V-\min(R,G,B)}\right), & V = B
     \end{cases}
\end{aligned}
\label{eq:hsv}
\end{equation}

\noindent dengan $R$, $G$, dan $B$ adalah nilai kanal merah, hijau, dan
biru yang telah dinormalisasi pada rentang 0 sampai 1. Nilai $H$
dinyatakan dalam derajat, sedangkan $S$ dan $V$ berada pada rentang 0
sampai 1.

%------------------------------------------------------------------
\section{Jaringan Saraf Tiruan}
%------------------------------------------------------------------

Jaringan saraf tiruan dengan algoritme perambatan balik
(\textit{backpropagation}) bekerja melalui dua tahap, yaitu perambatan
maju untuk menghitung keluaran dan perambatan balik untuk memperbaiki
bobot. Keluaran setiap neuron dihitung menggunakan fungsi aktivasi
sigmoid seperti pada \cref{eq:sigmoid}.

\begin{equation}
y_j = \frac{1}{1+e^{-\left(\sum_{i=1}^{n} w_{ij}x_i + b_j\right)}}
\label{eq:sigmoid}
\end{equation}

\noindent dengan $x_i$ adalah masukan ke-$i$, $w_{ij}$ adalah bobot dari
masukan ke-$i$ menuju neuron ke-$j$, dan $b_j$ adalah bias neuron ke-$j$.
Perbaikan bobot dilakukan dengan meminimumkan galat yang dinyatakan
pada \cref{eq:galat}.

\begin{equation}
E = \frac{1}{2}\sum_{j=1}^{m}\left(t_j-y_j\right)^2
\label{eq:galat}
\end{equation}

\noindent dengan $t_j$ adalah target keluaran dan $y_j$ adalah keluaran
jaringan. Bobot diperbarui sebesar $w_{ij} \leftarrow w_{ij} +
\eta\,\delta_j\,x_i$, dengan $\eta$ adalah laju pembelajaran dan
$\delta_j$ adalah gradien galat pada neuron ke-$j$.

%------------------------------------------------------------------
\section{Penelitian Terdahulu}
%------------------------------------------------------------------

Beberapa penelitian sebelumnya telah memanfaatkan citra digital untuk
menilai kesegaran ikan. Ringkasan penelitian tersebut disajikan pada
\cref{tab:terdahulu}.

\begin{table}[H]
\caption{Perbandingan Penelitian Terdahulu tentang Deteksi Kesegaran Ikan}
\label{tab:terdahulu}
\begin{spacing}{1}
\centering
\begin{tabularx}{\textwidth}{|>{\centering\arraybackslash}p{0.8cm}|X|>{\raggedright\arraybackslash}p{3.2cm}|>{\centering\arraybackslash}p{1.6cm}|}
\hline
\rowcolor{blue!12}
\textbf{No} & \textbf{Peneliti} & \textbf{Metode} & \textbf{Akurasi} \\
\hline
1 & \citet{prasetyo2021} & JST perambatan balik dengan ciri warna RGB & 88,00\% \\
\hline
2 & \citet{satria2004} & Analisis tekstur insang ikan & 84,50\% \\
\hline
3 & \citet{sudirman2010} & Klasifikasi warna mata ikan & 81,25\% \\
\hline
\end{tabularx}
\par\vspace{2pt}
\raggedright\noindent Sumber: hasil kajian pustaka (2024)
\end{spacing}
\end{table}

Berdasarkan \cref{tab:terdahulu}, sebagian besar penelitian menggunakan
ciri warna dan tekstur dengan akurasi di bawah 90\%. Penelitian ini
menggabungkan ciri warna pada ruang HSV dengan jaringan saraf tiruan
dua lapisan tersembunyi, yang diharapkan dapat meningkatkan akurasi
klasifikasi.
